% Section: Budget-Constrained Routing via Primal-Dual CBwK
% Include from paper/sections/results.tex

\subsection{Budget-Constrained Routing (Figure~\ref{fig:pareto_budget})}
\label{sec:budget_pacing}

\paragraph{Motivation.}
Production LLM deployments operate under cost constraints:
an engineering team may route millions of requests per day across
models spanning a $500\times$ price range
(e.g., \$0.00003/req for Llama-3.1-8B vs.\ \$0.015/req for Gemini-2.5-Pro)
and must keep average spend within an API budget.
The standard approach---a fixed cost-penalty weight $\lambda_s$ in the
bandit score---requires offline tuning to translate an abstract weight
into a dollar outcome, and must be re-tuned whenever the model portfolio
or its pricing changes.
We ask: \emph{can the router accept a dollar budget directly and
automatically achieve it, without sacrificing routing quality?}

\paragraph{Algorithm.}
We formulate the problem as a \emph{Contextual Bandit with Knapsacks}
(CBwK)~\cite{agrawal2014bandits} and solve it via a primal-dual scheme.
At each step $t$ the router selects model $a_t$ by maximising a
budget-augmented utility:
\begin{equation}
    a_t = \arg\max_{a \in \mathcal{A}_t} \Bigl[
        \underbrace{\hat\theta_a^\top x_t + \alpha \sqrt{x_t^\top A_a^{-1} x_t}}_{\text{LinUCB score}}
        \;-\; \underbrace{\lambda_s \,\tilde{c}_a}_{\text{static penalty}}
        \;-\; \underbrace{\lambda_t \,\tilde{c}_a}_{\text{dual penalty}}
    \Bigr]
    \label{eq:budget_ucb}
\end{equation}
where $\mathcal{A}_t \subseteq \mathcal{A}$ is the set of models
surviving a hard cost ceiling
$p_a \leq p_{\max}/(1 + \kappa\lambda_t)$,
$\tilde{c}_a \in [0,1]$ is the log-normalized unit cost,
and $\lambda_t$ is a Lagrange multiplier updated after observing the
realised cost $c_t$:
\begin{equation}
    \lambda_{t+1} = \min\!\Bigl(\bar\lambda,\;
    \max\bigl(0,\;\lambda_t + \eta\,(c_t / B - 1)\bigr)\Bigr)
    \label{eq:dual_update}
\end{equation}
with budget target $B$ (the user's desired average cost per request),
step size $\eta{=}0.05$, and cap $\bar\lambda{=}5$.
The normalisation $c_t / B$ makes $\eta$ portfolio-independent:
the same step size produces equivalent $\lambda_t$ trajectories regardless
of whether the cheapest model costs \$0.00003 or \$0.03.

The \textsc{BudgetPacer} operates in \emph{adaptive mode}, combining
two enforcement mechanisms:
\begin{itemize}[nosep,leftmargin=*]
  \item \textbf{Hard ceiling} (safety net):
    When $\lambda_t > 0$, the candidate set $\mathcal{A}_t$ excludes
    models whose blended unit price exceeds the dynamic ceiling,
    preventing catastrophic overspend on any single request.
  \item \textbf{Soft penalty} (optimizer):
    The dual penalty $\lambda_t\tilde{c}_a$ biases the UCB score toward
    cheaper models proportionally to their cost, allowing context-dependent
    routing within the surviving candidate set.
\end{itemize}

\noindent
This two-layer design mirrors production rate-limiting architectures:
the hard ceiling is a circuit breaker that guarantees worst-case
containment; the soft penalty is a fine-grained optimizer that
minimizes quality loss within the budget envelope.

\paragraph{Setup.}
We evaluate on the $K{=}3$ portfolio
(Llama-3.1-8B / Mistral-Large-2512 / Gemini-2.5-Pro)
under stationary conditions
(train: $n{=}8{,}374$; test: $n{=}1{,}824$; 32~NLP benchmarks;
5~seeds; forgetting factor $\text{ff}{=}1.0$).
The router is warm-started with offline priors
($n_{\text{eff}}{=}50$, $\alpha{=}1.0$, disjoint LinUCB).
Two conditions are compared:
\begin{enumerate}[nosep,leftmargin=*]
  \item \textbf{Static cost penalty:}
    Fixed $\lambda_s \in \{0, 0.05, 0.10, 0.20, 0.30, 0.50, 1.0\}$,
    no pacer.  Each value produces one (cost, reward) point.
  \item \textbf{BudgetPacer (adaptive):}
    Seven log-spaced targets from the cheapest model's empirical mean
    cost (\$0.000029/req) to the most expensive (\$0.0149/req),
    $\eta{=}0.05$, $\bar\lambda{=}5$.
\end{enumerate}

\paragraph{Pareto frontier (Figure~\ref{fig:pareto_budget}).}
Both methods trace a quality--cost Pareto curve on the test set.
At the budget-constrained end (cost $< \$0.001$/req), the pacer
achieves higher reward at lower cost:
at comparable spend (\$0.00008 vs.\ \$0.00012), the pacer
reaches mean reward $0.810$ vs.\ $0.797$ for the best static point
($\lambda_s{=}1.0$), a $+1.3$ percentage point improvement.
At moderate budgets (\$0.0007/req), the pacer achieves $0.863$ reward
at 20\% lower cost than the nearest static point
($\lambda_s{=}0.5$: $0.848$ at \$0.0008).

At the unconstrained end (cost $> \$0.005$), the static baseline
matches or slightly exceeds the pacer:
$\lambda_s{=}0$ achieves the quality ceiling ($0.928$) because it
imposes zero cost constraint.
The pacer at its loosest target ($B{=}\$0.015$) reaches $0.925$
at \$0.0076---$0.3$ points below the ceiling---because transient
$\lambda_t > 0$ episodes during the test phase create a small drag.
This convergence is expected: the pacer's value proposition is
\emph{budget compliance}, not Pareto dominance at every point.
When the budget is slack, the dual variable decays to zero and the
pacer gracefully becomes a no-op.

\paragraph{Budget compliance (Table~\ref{tab:budget_compliance}).}
The pacer's primary claim is that it \emph{hits the target}.
For the four tightest targets, budget utilisation ranges from
$0.96\times$ to $1.03\times$---effectively perfect compliance.
The two loosest targets ($B{=}\$0.005$ and $B{=}\$0.015$) show
under-utilisation ($0.86\times$ and $0.51\times$) because the
unconstrained router's natural spend (\$0.0082/req) is already below
these targets; the dual variable correctly decays to
$\lambda_t \approx 0$ rather than artificially inflating cost.
This asymmetry is by design: the pacer constrains spend
\emph{from above} but never forces the router to spend more than its
quality-optimal policy would.

\paragraph{Lambda dynamics (Figure~\ref{fig:lambda_convergence}).}
The dual-variable trajectory reveals the feedback loop.
At a tight target ($B{=}\$2.9{\times}10^{-5}$), $\lambda_t$
saturates near the cap ($\bar\lambda{=}5$) within the first 100
test steps and remains stable, producing a near-Llama-only policy.
At the moderate target ($B{=}\$6.6{\times}10^{-4}$), $\lambda_t$
oscillates around $0.3$--$0.5$: each Gemini selection causes a cost
spike that drives $\lambda_t$ up, which suppresses future Gemini
selections, allowing $\lambda_t$ to relax.
This oscillation is the hallmark of a primal-dual controller in
steady state: the system orbits the constraint boundary rather than
converging to a fixed point.
The median $\lambda_t$ (quartile $q_{50}{=}0.31$) is stable across
the trajectory, confirming convergence.

\paragraph{Model-mix dynamics (Figure~\ref{fig:model_mix_budget}).}
The stacked bar chart reveals a qualitative difference between the
two methods.
The static penalty produces a smooth, monotonic transition from
Gemini-dominant (51\% at $\lambda_s{=}0$) to Llama-dominant
(98\% at $\lambda_s{=}1.0$), with Mistral serving as an intermediate
buffer.
The pacer produces a \emph{sharper} and more principled transition:
at tight budgets the hard ceiling eliminates Gemini entirely (0\%
for $B \leq \$2.3{\times}10^{-4}$), and the soft penalty discovers a
Llama/Mistral blend (51\%/46\% at $B{=}\$6.6{\times}10^{-4}$) that
is unreachable by any single $\lambda_s$ value.
This blend exploits Mistral's favourable cost--quality ratio on prompts
where it matches Gemini's quality at $29\times$ lower cost---a
routing decision that the static penalty's uniform cost weighting
cannot express.

\paragraph{Production implications.}
\begin{enumerate}[leftmargin=*,noitemsep]
  \item \textbf{Interpretable interface.}
    The practitioner specifies ``keep average cost below
    \$X/request''---a quantity visible on any billing dashboard.
    The static penalty requires translating a dimensionless weight
    into a dollar outcome via offline simulation, and this mapping
    changes with every portfolio update (new model, price change,
    model deprecation).
    The pacer's dollar-denominated target is \emph{portable} across
    portfolios: the same $B{=}\$0.002$ target produces correct
    behaviour whether the portfolio has 2~models or 20, because the
    normalised dual update (Eq.~\ref{eq:dual_update}) adapts $\lambda_t$
    to the realised cost scale.

  \item \textbf{No offline tuning required.}
    The static baseline's 7-point Pareto curve requires 7~simulation
    runs to identify the operating point.
    The pacer achieves any point on its curve from a single parameter
    $B$, with no prior knowledge of the cost distribution.
    In a deployment where model pricing changes weekly (as is common
    with API providers), the pacer automatically adapts; the static
    penalty becomes stale.

  \item \textbf{Cost at scale.}
    Consider a deployment routing 100{,}000 requests/day.
    Moving from the unconstrained policy (\$0.0082/req, quality~0.928)
    to the pacer at $B{=}\$0.0019$/req saves
    $(\$0.0082 - \$0.0018) \times 100{,}000 = \$640$/day
    (\textbf{\$233{,}600/year})
    while maintaining quality~0.889---a 4.2\% quality reduction for a
    78\% cost reduction.
    The static baseline's nearest operating point
    ($\lambda_s{=}0.3$: \$0.0020/req, quality~0.896) delivers
    marginally higher quality but at 12\% higher cost (\$0.0020 vs.\
    \$0.0018), and offers no contractual guarantee that spend will
    remain at that level if the prompt distribution shifts.

  \item \textbf{Budget compliance as an SLA.}
    The $0.96$--$1.03\times$ utilisation range at tight targets
    means the pacer can underpin a cost SLA:
    ``average spend will not exceed $1.05 \times B$ over any
    rolling 1{,}000-request window.''
    The static penalty offers no such guarantee---its realised cost
    depends on the prompt distribution and can drift arbitrarily
    if traffic patterns change.

  \item \textbf{Graceful degradation.}
    When the budget is generous, the pacer does not distort routing:
    $\lambda_t \to 0$ and the quality-optimal policy is recovered.
    This means a team can set a conservative budget ceiling
    (e.g., $2\times$ expected spend) as a safety net without
    sacrificing quality in normal operation---the pacer only activates
    when actual spend approaches the ceiling.
\end{enumerate}
